% !TeX root = ../机器学习笔记.tex
% 由主文档通过 \input 载入；请编译主文档。
\setcounter{chapter}{0} % 保留原书章号。
\chapter{机器学习的任务与基本概念}
\label{chap:ml-tasks}

从面积预测房价、从邮件判断类别、让机器人选择动作，表面上是不同问题，却都涉及利用经验作出判断。本章先比较反馈机制，再区分数据、模型、目标和算法，建立后文的共同语言。

\emph{机器学习}（machine learning）是研究如何让计算机利用数据或交互经验，在给定任务和评价标准下改进表现的一类方法。其核心是从有限的经验中学习具有\emph{泛化能力}的模型，使模型能够对未见过的数据进行预测、发现数据中的潜在结构，或根据环境反馈选择合适的行动。因此，机器学习的目标不仅是拟合已有数据，更是利用从经验中获得的规律解决新的问题，典型任务包括分类、回归、聚类和序贯决策。

\section{监督学习、无监督学习与强化学习的区别}
\label{sec:learning-paradigms}

\emph{监督学习}（supervised learning）、\emph{无监督学习}（unsupervised learning）与\emph{强化学习}（reinforcement learning，RL）是机器学习的三种基本范式。理解它们的区别，需要同时考察三个问题：\emph{训练数据包含什么信息、学习过程依靠什么反馈、最终需要优化什么目标}。三者都可以使用神经网络等模型，因此不能仅凭模型结构判断一个任务属于哪种学习范式。

\begin{enumerate}
  \clubpenalty=10000
  \widowpenalty=10000
  \item \textbf{监督学习：从带标签的样本中学习预测规则}

  监督学习的训练集通常写为
  \[
    \mathcal{D}_{\mathrm{sup}}=\{(x_i,y_i)\}_{i=1}^{n},
  \]
  其中 $x_i$ 是输入随机变量 $X$ 的一个观测值，$y_i$ 是与之对应的\emph{标签}（label），即目标输出随机变量 $Y$ 的观测值。标签可以是离散类别，也可以是连续数值；前者对应\emph{分类}（classification），后者对应\emph{回归}（regression）。例如，根据已标注的邮件学习识别垃圾邮件，或根据房屋特征及其成交价格学习预测房价。

  设 $f_\theta$ 是参数为 $\theta$ 的预测模型，$\ell(\widehat y,y)$ 是衡量预测值与标签差异的\emph{损失函数}（loss function）。一种基本训练方式是\emph{经验风险最小化}（empirical risk minimization）：
  \begin{equation}
    \min_{\theta}\;\widehat{\mathcal{R}}_n(\theta),
    \qquad
    \widehat{\mathcal{R}}_n(\theta)
    =\frac{1}{n}\sum_{i=1}^{n}\ell\bigl(f_\theta(x_i),y_i\bigr).
  \end{equation}
  标签为每个样本提供直接的预测目标，但实际标签可能含有噪声。在训练样本独立同分布于 $P_{X,Y}$ 的常见设定下，最终关注的是模型对未见样本的\emph{泛化}（generalization）能力，即总体风险 $\mathcal{R}(\theta)=\mathbb{E}_{(X,Y)\sim P_{X,Y}}[\ell(f_\theta(X),Y)]$ 是否足够小，而不只是训练误差是否足够小。

  \item \textbf{无监督学习：从未标注的数据中提取结构}\nopagebreak

  无监督学习通常只有输入样本
  \[
    \mathcal{D}_{\mathrm{unsup}}=\{x_i\}_{i=1}^{n},
  \]
  没有针对当前任务给定的目标标签 $y_i$。学习任务可以是\emph{聚类}（clustering）、\emph{降维}（dimensionality reduction）或\emph{密度估计}（density estimation），分别关注样本的分组、低维表示或概率分布。例如，根据用户的行为特征发现相似群体，而不事先提供每位用户所属的群体类别。

  无监督学习并不意味着没有优化目标。以 $K$ 均值聚类（$K$-means clustering）为例，给定 $x_i\in\mathbb{R}^{d}$ 和簇数 $K$，可以求解
  \begin{equation}
    \min_{\mu_1,\ldots,\mu_K,\,c_1,\ldots,c_n}
    \sum_{i=1}^{n}\lVert x_i-\mu_{c_i}\rVert_2^2,
    \qquad c_i\in\{1,\ldots,K\},
  \end{equation}
  其中 $\mu_k\in\mathbb{R}^{d}$ 是第 $k$ 个簇的中心，$c_i$ 是算法为样本 $x_i$ 学到的簇编号。与监督学习的标签不同，$c_i$ 不是预先给定的目标答案。算法发现的结构取决于特征表示、距离度量和建模假设，也未必与人们期望的语义类别一致。

  \item \textbf{强化学习：从交互经验中学习序贯决策}

  强化学习研究\emph{智能体}（agent）如何在\emph{环境}（environment）中选择动作，以获得较高的长期收益。时刻 $t$，环境具有\emph{状态}（state）$S_t$，智能体接收\emph{观测}（observation）$O_t$，选择\emph{动作}（action）$A_t$；随后环境转移到 $S_{t+1}$，并给出\emph{奖励}（reward）$R_{t+1}$。状态描述环境的有关状况，观测是智能体实际获得的信息，两者一般不能等同；在完全可观测的设定下，通常简写为 $O_t=S_t$。

  智能体学习的是选择动作的\emph{策略}（policy）$\pi$。以完全可观测的离散问题为例，平稳\emph{随机策略}（stochastic policy）$\pi(a\mid s)$ 给出状态 $s$ 下选择动作 $a$ 的概率；\emph{确定性策略}（deterministic policy）则为每个状态指定一个动作。若只能部分观测环境，策略可能需要利用观测与动作的历史信息。

  对于奖励有界的无限时域任务，给定\emph{折扣因子}（discount factor）$0\leq\gamma<1$，一种常用目标为
  \begin{equation}
    G_t=\sum_{k=0}^{\infty}\gamma^k R_{t+k+1},
    \qquad J(\pi)=\mathbb{E}_{\pi}[G_0],
    \qquad \max_{\pi}\;J(\pi).
  \end{equation}
  这里 $G_t$ 称为\emph{回报}（return），是未来奖励的折扣累积；期望同时考虑给定的初始状态分布、策略和环境所产生的随机性。\emph{奖励是单步反馈，回报是跨时间累积的结果}，因此强化学习不等于逐步选择即时奖励最大的动作。

  与监督学习的标签相比，奖励通常只评价行为产生的结果，并不直接指出当前应当采取哪个动作。当前动作还可能改变未来状态，使奖励的影响延迟出现，形成\emph{信用分配}（credit assignment）问题，即如何判断先前哪些动作促成了后续结果。在线学习时，智能体还要权衡\emph{探索}（exploration）与\emph{利用}（exploitation）：探索通过尝试尚不熟悉的动作获得信息，利用则依据已有知识选择预期回报较高的动作。例如，机器人需要学习一系列移动动作，使整个搬运过程既高效又能避免碰撞。
\end{enumerate}

表~\ref{tab:learning-paradigms} 汇总了三类范式在典型设定下的区别。

\begin{table}[htbp]
  \centering
  \caption{三种学习范式的典型比较}
  \label{tab:learning-paradigms}
  \small
  \begin{tabularx}{\textwidth}{@{}L{1.7cm}YYY@{}}
    \toprule
    \textbf{比较维度} & \textbf{监督学习} & \textbf{无监督学习} & \textbf{强化学习} \\
    \midrule
    训练数据 & 输入与目标标签 $(x_i,y_i)$ & 未标注的输入 $x_i$ & 含观测、动作和奖励的交互轨迹 \\
    反馈依据 & 预测结果与给定标签之间的差异 & 数据本身及人为设定的结构准则 & 动作引起的奖励及其长期后果 \\
    学习目标 & 在新样本上准确预测 & 发现结构、表示或分布 & 学习使期望回报较大的策略 \\
    采样方式 & 通常使用预先给定的数据集 & 通常使用预先给定的数据集 & 在线交互时，策略会影响后续采样分布 \\
    典型任务 & 分类、回归 & 聚类、降维、密度估计 & 机器人控制、游戏决策、序贯资源调度 \\
    \bottomrule
  \end{tabularx}
\end{table}

这些区别描述的是\emph{任务的反馈机制与学习目标}，并非所有算法都严格采用同一种数据收集方式。例如，\emph{离线强化学习}（offline reinforcement learning）可以仅使用事先收集的交互数据学习策略，而不在训练时继续与环境交互；它仍然以优化序贯决策的回报为目标。

三类方法也可以组合使用：机器人可以用监督学习从带标签的图像中学习物体识别，用无监督学习提取传感器数据的结构，再用强化学习优化搬运策略。


\section{样本、特征、标签与预测任务}

一个\emph{样本}（sample）是所研究对象的一次观测，\emph{特征}（feature）是预测时可获得的描述，标签是希望预测的结果。例如，房屋输入 $x=(\text{面积},\text{楼龄})^\top$ 包含两个特征，输出 $y$ 是成交价格。随机向量 $X$ 描述可能出现的房屋，具体房屋的特征 $x_i$ 是其观测值；设计矩阵的一行对应一个样本，一列对应一个特征。

特征可以是数值、类别、文本或图像，但必须转化为模型能够处理的表示。类别编码为整数不代表类别间具有数值距离；缺失也不能无条件当作零。特征的可用时刻同样重要：预测成交价格时，只有成交之后才生成的费用记录通常不能用作输入。输入包含了预测时未知的信息，即使测试分数很高，也可能无法支持实际使用。

分类预测离散类别，回归预测数值；排序比较对象的相对次序，序列标注则为有顺序的多个位置赋予标签。相同数据可以服务于不同任务：用户分群没有预先给定的群体标签，而预测用户是否点击某物品具有行为标签。任务应由使用目的确定，不能由手中恰好有哪种算法反推。

\section{模型、学习目标与训练算法}

设一维输入为 $x$。选择 $f_\theta(x)=wx+b$，$\theta=(w,b)$，意味着只在直线组成的模型集合中寻找预测函数。选择平方损失意味着以预测偏差的平方评价单个样本；再选择梯度下降，才规定如何数值求解参数。三者分别回答“允许什么规律”“什么算好”“怎样找到它”。相同直线模型可使用平方或绝对损失，相同平方目标也可使用线性方程求解或迭代算法。

\emph{参数}（parameter）是由训练程序直接拟合的量，如 $w,b$；\emph{超参数}（hyperparameter）规定模型或训练程序，如多项式次数、正则化强度和学习率，通常利用验证数据选择。\emph{归纳偏置}（inductive bias）是学习方法在有限数据之外作出推广时采用的结构偏好，例如线性模型偏好线性关系，近邻方法偏好相近输入具有相近输出。这些偏好提供了学习的可能，也限定了方法的适用范围。

\section{训练、验证、测试与泛化}

训练集用于拟合参数，验证集用于比较模型和选择超参数，测试集用于在选择完成后评价整个流程。\emph{泛化}（generalization）关注未参与训练选择的新样本，训练误差则只描述已见样本。若算法在训练点逐一输出已记住的标签，而在其他位置恒输出零，那么它可以取得零训练误差，却没有学到可推广的关系。

测试误差也是有限样本上的估计，具有随机波动；它能否反映使用效果取决于测试样本是否代表目标分布以及是否独立于模型选择。反复查看测试分数并调整模型，相当于把测试数据用于验证，最终报告会受到选择偏差影响。增加模型容量常使训练误差下降，但可能增加对偶然噪声的依赖，形成\emph{过拟合}（overfitting）；更具体的理论见第\ref{chap:ml-generalization}章。

\section{学习范式的交叉与边界}

\emph{半监督学习}（semi-supervised learning）同时利用少量标注样本和较多未标注样本；\emph{自监督学习}（self-supervised learning）从数据自身构造预测任务，例如遮住部分输入再预测被遮住的内容。自监督方法仍有明确的训练目标，其特点是目标通常不依赖逐样本人工标注。预训练得到的表示还可以用于另一个有标签任务。

\emph{在线学习}（online learning）强调数据按时间到达并逐步更新模型，它不等于强化学习：逐封到达的邮件仍可构成监督分类，而强化学习的行动可能改变后续状态与反馈。神经网络是模型结构，监督或自监督是反馈组织方式，两种描述处于不同维度；因此“用了神经网络”不足以判定学习范式。

\section{例题、练习与学习检查}

\begin{example}[常数预测与学习目标]
训练响应为 $2,4,6$，只允许使用常数模型 $f_c(x)=c$。分别求平方损失和绝对损失下的最优预测。
\end{example}
\begin{solve}
平方损失为 $[(c-2)^2+(c-4)^2+(c-6)^2]/3$，导数为 $2(c-4)$，故最优常数是均值 $4$，平均损失为 $8/3$。绝对损失在中位数 $4$ 处最小，平均损失为 $4/3$。两种目标在本例给出相同参数，但损失数值不能直接比较为同一种评价；改变数据为 $2,4,60$ 后，两者的最优常数分别为 $22$ 和 $4$。
\end{solve}

习题：为预测房价、识别邮件和用户分群分别写出样本单位、输入、输出和可获得反馈；构造一个零训练误差但不能推广的函数；说明分类中永远输出训练多数类的模型为何仍值得作为基线。学习一种算法之前，应先能用这些问题清楚描述任务。





